\begin{table}[!htbp]
\centering
\caption{Sensitivity of addition returns}\label{tab:robustness}
\small
\setlength{\tabcolsep}{5pt}
\begin{tabular}{llrrl}
\toprule
Period & Specification & $N$ & CAR & 95\% interval \\
\midrule
2017--26 & Baseline & 169 & 2.15 & [1.01, 3.39] \\
2017--26 & Without MSCI overlap & 157 & 1.46 & [0.50, 2.47] \\
2017--26 & Without screened company news & 163 & 1.91 & [1.01, 2.82] \\
2017--26 & Winsorized 1st and 99th percentiles & 169 & 1.92 & [0.91, 3.05] \\
2017--26 & Without Sivers & 168 & 1.71 & [0.57, 2.96] \\
2022--26 & Baseline & 87 & 1.50 & [-0.10, 3.22] \\
2022--26 & Without MSCI overlap & 82 & 0.73 & [-0.77, 2.32] \\
2022--26 & Without screened company news & 84 & 1.30 & [0.12, 2.43] \\
2022--26 & Winsorized 1st and 99th percentiles & 87 & 1.04 & [-0.11, 2.25] \\
2022--26 & Without Sivers & 86 & 0.62 & [-0.85, 2.23] \\
\bottomrule
\end{tabular}
\begin{minipage}{\linewidth}\footnotesize\vspace{5pt}
Event-weighted CAR $[AD-1,ED+1]$ in percent, with pointwise restricted wild cluster bootstrap intervals by review. Each exclusion is applied separately. The company-news screen examines the 24 most extreme inherited events, rather than the full event universe, and is outcome-informed. MSCI exclusion changes the target population. Winsorization changes the outcome. These are sensitivity descriptions, not independent replications or specification selection rules. Sources: effect, msci\_overlap, news\_screen and extreme\_values aggregate CSV files.
\end{minipage}
\end{table}
